\subsection{Common Drain (CD, Source Follower)}\label{sec:amp_cd}

\subsubsection{Overview}
The input drives the gate, the output is taken at the source and the
drain is tied to the supply (AC ground). The output follows the input
shifted down by \(V_{GS}\) with a gain slightly below one. It has an
infinite input impedance and a low output impedance (\(\approx \frac{1}{g_m}\)):
a voltage buffer.

\subsubsection{Schematic}
\begin{circuitfig}
  \draw (5,0) node[nmos](M){};
  \node[right] at (5.4,0) {$M_1$};
  \draw (2.5,-2) node[ground]{} to[sinusoidal voltage source, l={$V_{in}$}] (2.5,0) -- (M.G);
  \draw (M.D) -- (5,2) node[vcc]{$V_{dd}$};
  \draw (M.S) -- (5,-1.4) node[circ]{} to[european resistor, l_={$G_L$}] (5,-3.4) node[ground]{};
  \draw (5,-1.4) -- (7.5,-1.4) node[circ]{} node[above]{$V_{out}$} to[capacitor, l={$C_L$}] (7.5,-3.4) node[ground]{};
\end{circuitfig}
\(G_L\) may be a resistor or the output conductance of a current source
(\(G_L = g_{o,sink}\)).

\subsubsection{Transfer Function}
\begin{enumerate}
  \item \textbf{Large Signal, Time Domain, DC Input Analysis (Q5)}: (\(K_{CLM} = 0\))
    \begin{align*}
      V_{out} G_L &= I_D = k_n (V_{in} - V_{out} - V_{th})^2 \\
      V_{in} - V_{out} - V_{th} &= \sqrt{\frac{V_{out} G_L}{k_n}} \\
      V_{out} &= V_{in} - V_{th} - \sqrt{\frac{V_{out} G_L}{k_n}} = V_{in} - V_{th} - V_{ov}
    \end{align*}
    With \(u = \sqrt{V_{out}}\): \(u^2 + \sqrt{\frac{G_L}{k_n}} u - (V_{in} - V_{th}) = 0\),
    \[
      V_{out} = \left(\frac{-\sqrt{\frac{G_L}{k_n}} + \sqrt{\frac{G_L}{k_n} + 4(V_{in} - V_{th})}}{2}\right)^2
    \]
    \(V_{th}\) itself rises with \(V_{out}\) through the body effect
    (\(V_{SB} = V_{out}\)). Saturation needs \(V_{dd} - V_{out} \geq V_{ov}\),
    i.e. \(V_{in} \leq V_{dd} + V_{th}\).
  \item \textbf{Small Signal, Frequency Domain, AC Input Analysis (Q4)}: \(V_d = 0\), \(V_b = 0\).\\
    Applying KCL at \(V_{out}\) (the drain current arrives at the source node):
    \begin{align*}
      V_{out}(G_L + sC_L) &= g_m (V_{in} - V_{out}) + g_{mb}(0 - V_{out}) + g_o (0 - V_{out}) \\
      V_{out}(G_L + g_m + g_{mb} + g_o + sC_L) &= g_m V_{in} \\
      H(s) = \frac{V_{out}}{V_{in}} &= \frac{g_m}{g_m + g_{mb} + g_o + G_L + sC_L}
    \end{align*}
\end{enumerate}

\subsubsection{Analysis}
\begin{enumerate}
  \item \textbf{DC Gain}:
    \[
      A_{V0} = \frac{g_m}{g_m + g_{mb} + g_o + G_L} \xrightarrow{G_L \to 0} \frac{g_m}{g_m + g_{mb} + g_o} \approx \frac{1}{1 + \frac{g_{mb}}{g_m}}
    \]
    The body effect limits the gain to about \(0.8\text{--}0.9\) unless the
    body is tied to the source (separate well).
  \item \textbf{Pole}: \(N(s)\) has no root, \(D(s) = 0\) at
    \[
      p_1 = \frac{g_m + g_{mb} + g_o + G_L}{C_L} \hspace{1cm} f_{p1} \approx \frac{g_m}{2\pi C_L}
    \]
    which is a very wide bandwidth (the GB of a CS stage driving the same
    \(C_L\)).
  \item \textbf{Impedances}:
    \paragraph{At Input:} \(I_G = 0\) so \(Z_{in} \to \infty\) at DC.
      \(C_{gs}\) is bootstrapped (both its ends move together), so it
      appears as \(C_{gs}(1 - A_{V0})\).
    \paragraph{At Output:} \(V_{in} = 0\), \(I_{test}\) into the source node
      (\(G_L\) removed):
      \begin{align*}
        I_{test} &= V_{out}(g_m + g_{mb} + g_o) \\
        Z_{out} &= \frac{1}{g_m + g_{mb} + g_o} \approx \frac{1}{g_m}
      \end{align*}
  \item \textbf{Level Shift}: \(V_{OUT} = V_{IN} - V_{th} - V_{ov}\); the stage
    costs \(V_{GS}\) of headroom.
\end{enumerate}

\subsubsection{Question}
A follower biased by a \(100\,\mu\)A sink (\(G_L = 2\,\mu\)S) has
\(V_{ov} = 0.2\,\)V, \(V_{th} = 0.5\,\)V, \(K_{CLM} = 0.1\,\)V\(^{-1}\),
\(g_{mb} = 0.2 g_m\), \(C_L = 5\,\)pF. Find \(A_{V0}\), \(Z_{out}\),
\(f_{p1}\) and \(V_{OUT}\) for \(V_{IN} = 1.2\,\)V.
\paragraph{Solution:}
\begin{align*}
  g_m &= 1\,mS \hspace{1cm} g_{mb} = 0.2\,mS \hspace{1cm} g_o = 10\,\mu S \\
  A_{V0} &= \frac{1\,m}{1\,m + 0.2\,m + 10\,\mu + 2\,\mu} = \frac{1}{1.212} = 0.825 \\
  Z_{out} &= \frac{1}{1.21\,mS} = 826\,\Omega \\
  f_{p1} &= \frac{1.212\,mS}{2\pi \cdot 5\,pF} = 38.6\,MHz \\
  V_{OUT} &= 1.2 - 0.5 - 0.2 = 0.5\,V \hspace{1cm} [\text{ignoring the body effect shift of } V_{th}]
\end{align*}
